\section{The GerryChain Baseline}
\label{sec:gerrychain-baseline}

\subsection{Comparison to MCMC Disruption}

A natural baseline for ``how much disruption is large?'' is the disruption
from a single \recom{} MCMC step, which randomises the plan by a small
but nonzero amount.
A full \recom{} chain of 10{,}000 steps has $d_{\rm Ham} \approx 1$ from
its start plan (approximately decorrelated).

We compare the reapportionment-induced boundary disruption to three
baselines:

\begin{table}[h]
\centering
\caption{Boundary disruption comparison across disruption types.
         $d_{\rm Ham}$ is the Hamming distance from the reference plan.}
\label{tab:disruption-comparison}
\begin{tabular}{lcc}
\toprule
Disruption type & $d_{\rm Ham}$ & Relative to reapportion \\
\midrule
Single \recom{} step (TX, $k=38$) & $\approx 0.026$ & $0.11\times$ \\
10 \recom{} steps (TX) & $\approx 0.22$ & $0.96\times$ \\
TX reapportionment ($38\to41$) & 0.23 & $1.0\times$ (baseline) \\
Seed change (worst case, GA) & 0.04 & $0.17\times$ \\
Political redistricting (typical) & 0.35--0.55 & $1.5$--$2.4\times$ \\
\bottomrule
\end{tabular}
\end{table}

The Texas reapportionment disruption ($d_{\rm Ham} = 0.23$) is equivalent
to approximately 10 \recom{} steps.
This is larger than the typical seed-change disruption (B.7) but
significantly smaller than political redistricting.

\subsection{Interpretation}

The finding that reapportionment disruption $\approx 10$ MCMC steps
has an important implication for statistical continuity:

\begin{enumerate}
\item \textbf{Algorithmic redistricting is more stable than MCMC.}
      A certified \recom{} ensemble (10{,}000+ steps) fully decorrelates
      from its starting plan.
      A reapportionment-driven redistricting changes only $\approx$23\% of
      tracts (for the most disruptive case, TX).
      The previous map is a useful starting point for the new map.
\item \textbf{Algorithmic redistricting is far more stable than political redistricting.}
      Political redistricting cycles (which change 35--55\% of tract
      assignments in contested states) produce $1.5$--$2.4\times$ more
      disruption than the worst algorithmic reapportionment.
\end{enumerate}

\subsection{Incumbent Protection Analysis}

One concern about reapportionment is that it disrupts incumbents by
moving them into different districts.
The Hamming distance measure gives an upper bound on incumbent disruption:
if $d_{\rm Ham} = 0.23$, at most 23\% of incumbents could be moved to
a different district (in practice fewer, as district assignments change
but incumbents can move with their home precinct).

For comparison, political redistricting cycles under the current system
often move 30--50\% of incumbents to new districts.
The \ar{} reapportionment is less disruptive to incumbents in the worst
case than typical political redistricting.

\subsection{GerryChain Baseline Definition}

The ``GerryChain baseline'' in the title of this section refers to using
the \recom{} step size ($d_{\rm Ham} \approx 1/k$ per step) as the unit
of measurement for boundary disruption.
This baseline is natural because:
\begin{enumerate}
\item It is the smallest unit of plan change in the MCMC framework.
\item It is computable for any state and any $k$.
\item It provides an apples-to-apples comparison between algorithmic
      stability and MCMC mixing.
\end{enumerate}

Under this baseline, the 2030 reapportionment for Texas is equivalent to
approximately $0.23 / 0.026 \approx 9$ \recom{} steps.
This is a concrete, policy-relevant characterisation: the 2030 Texas
reapportionment is a ``9-step'' disruption in MCMC terms.
